%% Отзыв научного руководителя (отдельный документ, не входит в main.tex).
%% Поля [...] заполняются научным руководителем / по данным диссертационного совета.
\documentclass[14pt,a4paper]{extarticle}

\usepackage[T2A]{fontenc}
\usepackage[utf8]{inputenc}
\usepackage[english,russian]{babel}

\usepackage{geometry}
\geometry{a4paper,top=2cm,bottom=2cm,left=3cm,right=1.5cm}

\usepackage{setspace}
\onehalfspacing
\usepackage{indentfirst}
\parindent=1.25cm
\sloppy
\pagestyle{empty}

\begin{document}

\begin{center}
    {\bfseries ОТЗЫВ НАУЧНОГО РУКОВОДИТЕЛЯ}

    \medskip
    {\bfseries Яниной Анастасии Олеговны}

    \medskip
    о работе Колодина Егора Ильича над диссертацией
    «Проблемы безопасности для больших языковых моделей в мультимодальном сценарии»,
    представленной на соискание учёной степени кандидата физико-математических наук
    по специальности 1.2.1 --- Искусственный интеллект и машинное обучение
\end{center}

\bigskip

Колодин Егор Ильич выполнил диссертационное исследование под научным руководством автора настоящего отзыва. За время работы над диссертацией Е.\,И.~Колодин зарекомендовал себя как квалифицированный специалист в области обработки естественного языка, информационного поиска и обеспечения безопасности больших языковых моделей. Он продемонстрировал способность самостоятельно ставить и решать сложные научно-технические задачи, проводить теоретические и экспериментальные исследования, анализировать и обобщать полученные результаты, а также уверенно ориентироваться в современной научной литературе по тематике работы.

Тема диссертации актуальна. Промышленное развёртывание русскоязычных и мультимодальных больших языковых моделей требует одновременно качественного информационного поиска --- для заземления ответов на достоверных источниках --- и надёжных механизмов обеспечения безопасности при сохранении полезности модели. Оба направления для русского языка проработаны существенно слабее, чем для английского, а их интеграция в единый защитный контур ранее систематически не рассматривалась. Ключевая гипотеза работы --- о том, что безопасность мультимодальных моделей должна обеспечиваться не только фильтрацией входов и выходов, но и контролем информационного пространства, из которого модель получает внешние знания, --- является научно содержательной и практически значимой.

Научная новизна работы определяется рядом полученных результатов. Предложен иерархический трёхэтапный конвейер инструктивного дообучения декодерной языковой модели для построения русскоязычных текстовых эмбеддингов (GigaEmbeddings); модель занимает первое место в публичном рейтинге MTEB(rus) при более чем вдвое меньшем числе параметров по сравнению с сильнейшими базовыми моделями. Разработан кросс-энкодер переранжирования на основе компактных декодерных моделей с заменой каузального механизма внимания на двунаправленный, агрегацией mean pooling и выделенной классификационной головой. Предложен двухэтапный конвейер обеспечения безопасности мультимодальных моделей, сочетающий независимый от защищаемой модели потоковый классификатор обменов с безопасной моделью, дообученной с подкреплением с гибридной наградой за безопасность. Впервые поисковый стек на основе эмбеддингов и переранжирования интегрирован в конвейер безопасности в виде инструмента безопасного поиска по курируемому индексу, а защитные механизмы распространены на визуальную модальность.

Практическую ценность работы представляют разработанные модели и методы, применимые в промышленных русскоязычных системах. Модель эмбеддингов и базовые модели семейства выпущены в открытый доступ, а предложенный конвейер обеспечения безопасности рассчитан на развёртывание с низкими накладными расходами и независимое обновление защитного слоя относительно защищаемой модели.

Основные результаты диссертационной работы опубликованы в двух рецензируемых статьях в трудах 63-й конференции Ассоциации компьютерной лингвистики (ACL~2025):
\begin{enumerate}
    \item GigaChat team. GigaChat Family: Efficient Russian Language Modeling Through Mixture of Experts Architecture // Proceedings of the 63rd Annual Meeting of the Association for Computational Linguistics (Volume~3: System Demonstrations). --- Association for Computational Linguistics, 2025. --- P.~93--106.
    \item Kolodin~E., Khomich~D., Savushkin~N., Ianina~A., Minkin~F. GigaEmbeddings --- Efficient Russian Language Embedding Model // Proceedings of the 10th Workshop on Slavic Natural Language Processing (Slavic~NLP~2025). --- Association for Computational Linguistics, 2025. --- P.~17--24.
\end{enumerate}
В статье~[2], посвящённой ключевому результату работы --- модели GigaEmbeddings, --- Е.\,И.~Колодин является первым автором. Результаты прошли апробацию на конференции ACL~2025 (программа System Demonstrations) и на воркшопе Slavic NLP~2025; модели GigaEmbeddings и GigaChat выпущены в открытый доступ. Кроме того, Е.\,И.~Колодин осуществлял научное со-руководство магистерской диссертацией, посвящённой реранкерам для русскоязычного семантического поиска, результаты которой использованы в настоящей работе с явным указанием личного вклада.

Основные результаты, выносимые на защиту, получены лично соискателем; вклад соискателя в коллективно выполненные работы (семейство моделей GigaChat) и в со-руководимые исследования чётко разграничен во введении к диссертации. В процессе подготовки диссертации Е.\,И.~Колодин проявил умение работать с обширной базой научных источников, настойчивость в достижении поставленных целей и хорошую ориентацию в специфике исследуемой предметной области.

Колодина Егора Ильича можно охарактеризовать как сформировавшегося научного работника и исследователя, способного самостоятельно ставить и решать сложные научно-технические задачи.

Тематика работы соответствует направлениям исследований паспорта научной специальности 1.2.1 «Искусственный интеллект и машинное обучение», в том числе разработке методов и алгоритмов для обработки и анализа текстов на естественном языке и изображений (п.~4), методам и технологиям поиска, приобретения и использования знаний в системах искусственного интеллекта (п.~5), а также исследованиям в области «доверенных» систем класса искусственного интеллекта, включая вопросы их надёжности и устойчивости (п.~12).

Диссертационная работа Колодина Егора Ильича является завершённым научно-квалификационным исследованием, посвящённым актуальной научно-технической задаче, и соответствует требованиям, предъявляемым к диссертациям на соискание учёной степени кандидата физико-математических наук, а также паспорту специальности 1.2.1 «Искусственный интеллект и машинное обучение». Рекомендую диссертацию «Проблемы безопасности для больших языковых моделей в мультимодальном сценарии» к защите в диссертационном совете.

\vspace{2.5em}

\noindent
\begin{minipage}{0.55\textwidth}
Научный руководитель,\\
кандидат физико-математических наук,\\
научный сотрудник лаборатории фундаментальных исследований искусственного интеллекта МФТИ
\end{minipage}
\hfill
\begin{minipage}{0.4\textwidth}
\raggedleft
\rule{4cm}{0.4pt}\\
Янина А.\,О.

\medskip
«\rule{1cm}{0.4pt}» \rule{3cm}{0.4pt} 2026~г.
\end{minipage}

\end{document}
